% Generated draft: requires mathematical and source review.
\documentclass[11pt]{article}
\usepackage[margin=25mm]{geometry}
\usepackage{fontspec}
\setmainfont{Latin Modern Roman}
\usepackage{amsmath,amssymb,mathtools,unicode-math}
\setmathfont{Latin Modern Math}
\usepackage{enumitem,needspace,xurl,hyperref}
\setlength{\parindent}{0pt}
\setlength{\parskip}{5pt}
\setlength{\emergencystretch}{3em}
\newcommand{\N}{\mathbb{N}}
\newcommand{\Z}{\mathbb{Z}}
\newcommand{\Q}{\mathbb{Q}}
\newcommand{\R}{\mathbb{R}}
\newcommand{\C}{\mathbb{C}}
\newcommand{\F}{\mathbb{F}}
\newcommand{\A}{\mathbb{A}}
\newcommand{\PP}{\mathbb{P}}
\newcommand{\E}{\mathbb{E}}
\newcommand{\eps}{\varepsilon}
\newcommand{\dd}{\,\mathrm{d}}
\begin{document}
\section*{UnsolvedMath ID: 30000489}
\textbf{Source problem code:} catalogue entry
\subsection{Definitions and mathematical statement}
Fix a prime \(p\). In the usual generalized Greenberg setting, for a number field \(F\), let \(\widetilde F\) be the compositum of all \(\Z_p\)-extensions of \(F\) in a fixed algebraic closure, and put
\[
\Gamma=\operatorname{Gal}(\widetilde F/F),\qquad
\widetilde\Lambda=\Z_p[[\Gamma]].
\]
The group \(\Gamma\) is a finite-rank \(\Z_p\)-module in this convention; no Leopoldt-based formula for its rank is assumed. Let \(L/\widetilde F\) be the maximal unramified abelian pro-\(p\) extension in the intended ramification convention, and write
\[
X_{\widetilde F}=\operatorname{Gal}(L/\widetilde F),
\]
with the \(\widetilde\Lambda\)-module action induced by conjugation over \(F\).

For a finitely generated module over this commutative Iwasawa algebra, pseudo-nullity means that localization is zero at every prime ideal of height at most one; equivalently its support has codimension at least two. This does not mean that the module itself is zero. The source's Greenberg conjecture is
\[
\forall F\text{ a number field},\qquad
X_{\widetilde F}\text{ is pseudo-null over }\widetilde\Lambda.
\]
The statement does not restrict \(F\) to be totally real. The prime range, precise definition of \(\widetilde F\), unramified versus split-at-selected-primes convention for \(X\), archimedean conditions at \(p=2\), and finite-generation assumptions are not expanded in the extract; the definitions above are the standard interpretation requiring source matching. No noncommutative Iwasawa algebra is substituted for this compositum setting.

The supplied literature assessment checked 2026-08-21 reports a cyclotomic-field case for \(p<1000\) and additional explicit nonabelian examples, while saying the arbitrary totally real case remains open. It does not specify the exact cyclotomic fields, prime exclusions, hypotheses, or what nonabelian means in those examples. These partial descriptions are not promoted to a theorem for every field in the displayed universal conjecture.
\subsection{Short English statement}
For the usual compositum of all \(\Z_p\)-extensions of a number field, ask whether its unramified abelian pro-\(p\) Galois module has support of codimension at least two over the Iwasawa algebra. The source asks about every number field, not only totally real ones.
\subsection{Sources}
\begin{itemize}
\item Local UnsolvedMath registry, permanent ID 30000489, catalogue entry, \emph{Greenberg's Pseudo-Nullity Conjecture}; attribution to Greenberg is explicit in the original statement.
\item \emph{Pro-\(p\) Extensions of Global Fields and pro-\(p\) Groups} (2007), Oberwolfach Reports, pp. 1463--1536, DOI \url{https://doi.org/10.4171/owr/2006/25}.
\item Supplied reference ``Sharifi primary paper'': \url{https://arxiv.org/abs/0801.1360}.
\item Supplied reference ``Jaulent on Greenberg conjecture'': \url{https://arxiv.org/abs/1612.00718}.
\item Supplied literature triage checked 2026-08-21; exact special-case statements and notation definitions are not reproduced.
\end{itemize}
\end{document}
